% Generated draft: requires mathematical and source review.
\documentclass[11pt]{article}
\usepackage[margin=25mm]{geometry}
\usepackage{fontspec}
\setmainfont{Latin Modern Roman}
\usepackage{amsmath,amssymb,mathtools,unicode-math}
\setmathfont{Latin Modern Math}
\usepackage{enumitem,needspace,xurl,hyperref}
\setlength{\parindent}{0pt}
\setlength{\parskip}{5pt}
\setlength{\emergencystretch}{3em}
\newcommand{\N}{\mathbb{N}}
\newcommand{\Z}{\mathbb{Z}}
\newcommand{\Q}{\mathbb{Q}}
\newcommand{\R}{\mathbb{R}}
\newcommand{\C}{\mathbb{C}}
\newcommand{\F}{\mathbb{F}}
\newcommand{\A}{\mathbb{A}}
\newcommand{\PP}{\mathbb{P}}
\newcommand{\E}{\mathbb{E}}
\newcommand{\eps}{\varepsilon}
\newcommand{\dd}{\,\mathrm{d}}
\begin{document}
\section*{UnsolvedMath ID: 2939}
\textbf{Source problem code:} catalogue entry
\subsection{Definitions and mathematical statement}
In the usual setting of Seiberg--Witten invariants, let $X$ be a closed connected oriented smooth $4$-manifold with $\pi_1(X)=1$, with the additional invariant/chamber choices required below. Let $Q_X$ be the integral intersection form on $H_2(X;\Z)$, and let $b_2(X)$ and $b_2^\pm(X)$ be its rank and positive/negative indices. The form is indefinite when both indices are positive.

For the chosen Seiberg--Witten invariant, let $\mathcal B(X)=\{c_1(\mathfrak s):\operatorname{SW}_X(\mathfrak s)\ne0\}\subseteq H^2(X;\Z)$, where $\mathfrak s$ is a spin-$c$ structure. Define
\[\operatorname{br}(X)=\dim_{\Q}\operatorname{Span}_{\Q}\{b\otimes1:b\in\mathcal B(X)\}\subseteq H^2(X;\Q),\]
with value zero if there are no basic classes. This is the rank-of-span notion in the source, not the number of basic classes. Closedness, orientation, invariant version, and homology-orientation conventions are not explicit in the stored statement. In particular $b_2^+=1$ is included in its odd-$b_2^+$ range and requires chamber choices; none is silently supplied or removed by restricting to $b_2^+>1$.

The three questions are:
\begin{itemize}
\item (a) Find an upper bound for $\operatorname{br}(X)$ in terms of topological invariants of $X$. The tautological bound $0\leq\operatorname{br}(X)\leq b_2(X)$ follows from its definition; to be a nontrivial open request this must seek additional strength, but the source gives no exact required improvement beyond part (b).
\item (b) Is $\operatorname{br}(X)\leq b_2^+(X)$ for every simply connected $X$ in the intended scope with $b_2^+(X)$ odd?
\item (c) Does there exist an indefinite simply connected smooth $X$ with $b_2(X)\geq3$ such that the induced diffeomorphism action on $Q_X$ has finite image, or even trivial image?
\end{itemize}
For oriented $X$, an orientation-preserving diffeomorphism acts on $H_2(X;\Z)$ by an integral isometry, giving $ev_*:\operatorname{Diff}^+(X)\to\operatorname{Aut}(Q_X)$. The source writes $\operatorname{Diff}(X)$ without a superscript; whether it means this group must be checked, since an orientation-reversing map generally reverses the intersection pairing. Trivial image means every allowed diffeomorphism acts as the identity, a stronger requirement than finite image.

The source describes Fintushel--Stern knot surgery on independent disjoint tori under nonvanishing Seiberg--Witten hypotheses as producing independent basic-class directions, and cites Gompf--Mrowka examples with $\operatorname{br}(X)=b_2^+(X)$. Those examples show the proposed bound's possible sharpness rather than prove it universally. The notes explain that diffeomorphisms permute basic classes up to sign; a basic-class span of full rank $b_2(X)$ would force the isometry image to be finite. In the indefinite setting $b_2^+(X)<b_2(X)$, so the cited rank-$b_2^+$ examples alone do not establish that full-rank mechanism or part (c). The indefinite/rank-at-least-three restriction is intended to avoid a trivially finite ambient form automorphism group.

The supplied triage checked 2026-08-17 labels the programme partially solved / PARTIAL-PROGRESS and reports neither the intended stronger general bound nor the required finite/trivial-image example verified. The registry is open. The literal topological scope, chamber, action-group, and nontrivial-strength conventions remain explicit formulation issues.
\subsection{Short English statement}
How large can the span of Seiberg--Witten basic classes be, especially compared with $b_2^+$, and can their constraints make the diffeomorphism action on an indefinite intersection form finite or trivial? The source reports sharp-rank examples but no universal inequality or required action example, with several scope and chamber conventions missing.
\subsection{Sources}
\begin{itemize}
\item Local UnsolvedMath registry, permanent ID 2939, source code catalogue entry, ``Kirby Problem 4.63''; supplied basic-rank definition, three questions, constructions, and literature triage.
\item Baykur, Kirby, and Ruberman, \emph{K3 -- A New Problem List in Low-Dimensional Topology}, Problem 4.63 (2026): \url{https://bpb-us-e2.wpmucdn.com/websites.umass.edu/dist/b/22144/files/2026/04/K3-problem-list-watermarked.pdf}. Supplied maintained source for the questions and construction context.
\item [FS98] Ronald Fintushel and Ronald J. Stern, ``Knots, links, and 4-manifolds,'' \emph{Inventiones Mathematicae} 134(2) (1998), 363--400: \url{https://doi.org/10.1007/s002220050268}. Supplied knot-surgery source.
\item [GM93] Robert Gompf and Tomasz Mrowka, ``Irreducible 4-manifolds need not be complex,'' \emph{Annals of Mathematics} 138 (1993), 61--111: \url{https://doi.org/10.2307/2946635}. Supplied construction source for rank-$b_2^+$ examples.
\item Supplied literature triage checked 2026-08-17, PARTIAL-PROGRESS; it distinguishes examples and sharpness scenarios from a universal bound or part-(c) solution, flags the merged-word/split-inequality extraction defects, and cautions that dated triage does not exclude later or unindexed results.
\end{itemize}
\end{document}
